% appendix.tex: VERACISM Test Results Appendices
% NOTE: \appendix is intentionally omitted -- it conflicts with titlesec in the
% report class used by is295.cls and causes an infinite loop compile hang.
% Section titles carry the "APPENDIX X" prefix explicitly instead.

% ============================================================
\section{APPENDIX A: INTEGRATION TEST SUITE RESULTS}
\label{appendix:integration-tests}
% ============================================================

The final results for all assertions contained within each of the following tables are the results of the VERACISM API v1 Integration Test run on 30 March 2026, using the Newman CLI runner against the Veracism staging environment
(\texttt{https://sb-api.ismanila.org/Veracism}). No failures or errors in the 93 assertions made for 22 test suites.

\subsection*{A.1: Scan Upload (\texttt{POST /api/v1/scan-upload})}

\begin{table}[H]
  \centering
  \smalltable
  \footnotesize
  {%
  \setlength{\tabcolsep}{5pt}
  \renewcommand{\arraystretch}{1.2}
  \begin{tabular}{|>{\raggedright\arraybackslash}p{0.07\textwidth}|
                  >{\raggedright\arraybackslash}p{0.52\textwidth}|
                  >{\centering\arraybackslash}p{0.12\textwidth}|
                  >{\centering\arraybackslash}p{0.13\textwidth}|}
    \hline
    \textbf{TC ID} & \textbf{Assertion} & \textbf{Result} & \textbf{Time (s)} \\
    \hline\hline
    TC-01 & Status 200; SHA-256 is 64-char hex; threatLevel valid; reportId is a number; results array present; requestId (GUID or 32-hex); response time \textless{}30 s & \textbf{PASS} & 20.446 \\
    \hline
    TC-02 & Status 200; studentId is non-null when studentNumber provided; response time \textless{}30 s & \textbf{PASS} & 17.321 \\
    \hline
    TC-03 & Status 200; versionNumber $\geq$ 2; documentLabel present; response time \textless{}30 s & \textbf{PASS} & 17.392 \\
    \hline
    TC-04 & Status 400; error is ``no files''; response time \textless{}30 s & \textbf{PASS} & 0.054 \\
    \hline
    TC-05 & Status 401 or 403 for invalid API key; response time \textless{}30 s & \textbf{PASS} & 0.134 \\
    \hline
    TC-06 & Status 401 or 403 when no \texttt{X-Api-Key} header; response time \textless{}30 s & \textbf{PASS} & 0.050 \\
    \hline
    TC-07 & Status 200; results array present; SHA-256 present for TXT upload; response time \textless{}30 s & \textbf{PASS} & 16.888 \\
    \hline
    TC-07A & Status 200; YARA match detected; webshell signature identified; threat level not SAFE; Lighthouse upload blocked; SHA-256 is 64-char hex; response time \textless{}30 s & \textbf{PASS} & 4.734 \\
    \hline
  \end{tabular}
  }%
  \captionsetup{type=table}
  \caption{Integration test results: Scan Upload group (37 assertions, 0 failures, 77.969 s)}
  \label{tab:int-scan-upload}
\end{table}

\subsection*{A.2: Get Report by CID (\texttt{GET /api/v1/reports/\{cid\}})}

\begin{table}[H]
  \centering
  \smalltable
  \footnotesize
  {%
  \setlength{\tabcolsep}{5pt}
  \renewcommand{\arraystretch}{1.2}
  \begin{tabular}{|>{\raggedright\arraybackslash}p{0.07\textwidth}|
                  >{\raggedright\arraybackslash}p{0.52\textwidth}|
                  >{\centering\arraybackslash}p{0.12\textwidth}|
                  >{\centering\arraybackslash}p{0.13\textwidth}|}
    \hline
    \textbf{TC ID} & \textbf{Assertion} & \textbf{Result} & \textbf{Time (s)} \\
    \hline\hline
    TC-08 & Status 200; CID matches; SHA-256 is 64-char hex; filename non-empty; status is ``complete''; verificationUrl present; threatLevel valid; uploadedAt is a date; response time \textless{}30 s & \textbf{PASS} & 0.147 \\
    \hline
    TC-09 & Status 404; error is ``Report not found''; response time \textless{}30 s & \textbf{PASS} & 0.054 \\
    \hline
    TC-10 & Status 401 or 403 with no API key; response time \textless{}30 s & \textbf{PASS} & 0.048 \\
    \hline
    TC-11 & Status 404 for non-owned CID (tenant enumeration hidden); response is not 403; response time \textless{}30 s & \textbf{PASS} & 0.053 \\
    \hline
  \end{tabular}
  }%
  \captionsetup{type=table}
  \caption{The integration tests returned 17 assertions, 0 failures, and 0.307 s overall.The total time for the integration tests is 0.307 s, with 0 failures, 17 assertions, and a total of 0 reports.}
  \label{tab:int-report-cid}
\end{table}

\subsection*{A.3: Public Verification, Download, and Storage Status}

\begin{table}[H]
  \centering
  \smalltable
  \footnotesize
  {%
  \setlength{\tabcolsep}{5pt}
  \renewcommand{\arraystretch}{1.2}
  \begin{tabular}{|>{\raggedright\arraybackslash}p{0.07\textwidth}|
                  >{\raggedright\arraybackslash}p{0.10\textwidth}|
                  >{\raggedright\arraybackslash}p{0.42\textwidth}|
                  >{\centering\arraybackslash}p{0.12\textwidth}|
                  >{\centering\arraybackslash}p{0.10\textwidth}|}
    \hline
    \textbf{TC ID} & \textbf{Group} & \textbf{Key Assertions} & \textbf{Result} & \textbf{Time (s)} \\
    \hline\hline
    TC-12 & Verify & verified is true; CID matches; SHA-256 64-char hex; uploadedAt present; threatLevel valid & \textbf{PASS} & 0.066 \\
    \hline
    TC-13 & Verify & Status 200; verified is false; message present (unknown CID) & \textbf{PASS} & 0.055 \\
    \hline
    TC-14 & Verify & Status 200; verified is false for garbage CID; message present & \textbf{PASS} & 0.053 \\
    \hline
    TC-15 & Verify & Status 200 without API key (public endpoint); verified is true & \textbf{PASS} & 0.056 \\
    \hline
    TC-16 & Download & Status 200; Content-Disposition: attachment; Content-Type present; body non-empty & \textbf{PASS} & 5.841 \\
    \hline
    TC-17 & Download & Status 404; error is ``Report not found'' & \textbf{PASS} & 0.054 \\
    \hline
    TC-18 & Download & Status 401 or 403 when no API key & \textbf{PASS} & 0.119 \\
    \hline
    TC-19 & Storage & Status 200; status is ``pinned'' or ``processing''; accessible is boolean; size is a number & \textbf{PASS} & 0.061 \\
    \hline
    TC-20 & Storage & Status 404; error is ``Report not found'' & \textbf{PASS} & 0.053 \\
    \hline
    TC-21 & Storage & Status 401 or 403 when no API key & \textbf{PASS} & 0.049 \\
    \hline
  \end{tabular}
  }%
  \captionsetup{type=table}
  \caption{Integration test results: Verify, Download, and Storage Status groups (41 assertions, 0 failures)}
  \label{tab:int-verify-download-storage}
\end{table}

\subsection*{A.4: Integration Suite Overall Summary}

\begin{table}[H]
  \centering
  \smalltable
  \footnotesize
  {%
  \setlength{\tabcolsep}{6pt}
  \renewcommand{\arraystretch}{1.3}
  \begin{tabular}{|>{\raggedright\arraybackslash}p{0.28\textwidth}|
                  >{\centering\arraybackslash}p{0.12\textwidth}|
                  >{\centering\arraybackslash}p{0.14\textwidth}|
                  >{\centering\arraybackslash}p{0.08\textwidth}|
                  >{\centering\arraybackslash}p{0.08\textwidth}|
                  >{\centering\arraybackslash}p{0.14\textwidth}|}
    \hline
    \textbf{Group} & \textbf{Test cases} & \textbf{Assertions} & \textbf{Pass} & \textbf{Fail} & \textbf{Duration (s)} \\
    \hline\hline
    Scan Upload & 8 & 37 & 37 & 0 & 77.969 \\
    \hline
    Get Report by CID & 4 & 17 & 17 & 0 & 0.307 \\
    \hline
    Verify (Public) & 4 & 18 & 18 & 0 & 0.230 \\
    \hline
    Download & 3 & 12 & 12 & 0 & 6.014 \\
    \hline
    Storage Status & 3 & 11 & 11 & 0 & 0.163 \\
    \hline
    \textbf{TOTAL} & \textbf{22} & \textbf{93} & \textbf{93} & \textbf{0} & \textbf{83.728} \\
    \hline
  \end{tabular}
  }%
  \captionsetup{type=table}
  \caption{Integration test suite overall summary
}
  \label{tab:int-overall}
\end{table}

% ============================================================
\clearpage
\section{APPENDIX B: FUNCTIONAL TEST SUITE RESULTS}
\label{appendix:functional-tests}
% ============================================================


The results of the entire VERACISM Functional Test Suite were run on 31 March 2026 and are presented in the following tables.  The 15 automated API test cases take place through automation with Bearer-token authentication, executed through Newman; 11 verifications are made though Code Review; 4 are Observable checks.  0 failures and all 89 automated assertions passed.

\subsection*{B.1: Role-Based Functional Test Results}

The 30 functional test cases are included below, each listed in a separate table, one for each actor role, for ease of readability.

%: B.1a System Administrator ---
\begin{table}[H]
  \centering
  \smalltable
  \footnotesize
  {%
  \setlength{\tabcolsep}{4pt}
  \renewcommand{\arraystretch}{1.3}
  \begin{tabular}{|>{\raggedright\arraybackslash}p{0.12\textwidth}|
                  >{\raggedright\arraybackslash}p{0.27\textwidth}|
                  >{\raggedright\arraybackslash}p{0.11\textwidth}|
                  >{\centering\arraybackslash}p{0.10\textwidth}|
                  >{\raggedright\arraybackslash}p{0.26\textwidth}|}
    \hline
    \textbf{TC ID} & \textbf{Test Case} & \textbf{Method} & \textbf{Result} & \textbf{Key Evidence} \\
    \hline\hline
    TC-ADM-001 & Approve integration request & Code Review & \textbf{PASSED} & Admin tenant management pages; \texttt{TenantsController} CRUD; \texttt{is\_active} flag \\
    \hline
    TC-ADM-002 & Manage users and roles & Code Review & \textbf{PASSED} & \texttt{/admin/users} UI; \texttt{UsersController}; Auth0 role sync via \texttt{ValidateUserExistsFilter} \\
    \hline
    TC-ADM-003 & Update keys and critical config & Code Review & \textbf{PASSED} & API Keys settings page; \texttt{ApiClient} table; ClamAV/YARA config via \texttt{appsettings.json} \\
    \hline
    TC-ADM-004 & View audit logs and export & Code Review & \textbf{PASSED} & \texttt{GET /audit/export/csv} and \texttt{/pdf}; actor, action, and tenantId filtering; 10,000-record export cap \\
    \hline
    TC-ADM-005 & View health and alerts & Code Review & \textbf{PASSED} & Admin dashboard health widgets; ClamAV/Lighthouse/DB status indicators \\
    \hline
  \end{tabular}
  }%
  \captionsetup{type=table}
  \caption{Functional tests: System Administrator (TC-ADM), 5 test cases}
  \label{tab:func-adm}
\end{table}

%: B.1b IT Office Staff ---
\begin{table}[H]
  \centering
  \smalltable
  \footnotesize
  {%
  \setlength{\tabcolsep}{4pt}
  \renewcommand{\arraystretch}{1.3}
  \begin{tabular}{|>{\raggedright\arraybackslash}p{0.11\textwidth}|
                  >{\raggedright\arraybackslash}p{0.30\textwidth}|
                  >{\raggedright\arraybackslash}p{0.15\textwidth}|
                  >{\centering\arraybackslash}p{0.10\textwidth}|
                  >{\raggedright\arraybackslash}p{0.22\textwidth}|}
    \hline
    \textbf{TC ID} & \textbf{Test Case} & \textbf{Method} & \textbf{Result} & \textbf{Key Evidence} \\
    \hline\hline
    TC-ITO-001 & API upload integration & Automated API & \textbf{PASSED} & \texttt{POST /api/v1/scan-upload} returns requestId, sha256, reportId, cid; all 7 assertions pass \\
    \hline
    TC-ITO-002 & Unauthorized upload rejection & Automated API & \textbf{PASSED} & Invalid key $\rightarrow$ 401/403; missing header $\rightarrow$ 401/403; no file stored \\
    \hline
    TC-ITO-003 & Metadata consistency & Automated API & \textbf{PASSED} & \texttt{GET /api/v1/reports/\{cid\}} CID, SHA-256, filename, uploadedAt, status all consistent \\
    \hline
    TC-ITO-004 & Lighthouse anchoring & Automated API & \textbf{PASSED} & CID non-empty; status is ``pinned'' or ``processing''; accessible: true; file size positive \\
    \hline
    TC-ITO-005 & Temporary Lighthouse outage retry & Code Review & \textbf{FAILED} & No Polly retry policy; single HTTP POST with no retry/circuit breaker; future enhancement \\
    \hline
    TC-ITO-006 & Scoped audit access & Code Review & \textbf{PARTIAL} & RBAC enforced; tenant isolation on Report queries; dedicated audit log UI not yet built \\
    \hline
  \end{tabular}
  }%
  \captionsetup{type=table}
  \caption{Functional tests: IT Office Staff (TC-ITO), 6 test cases}
  \label{tab:func-ito}
\end{table}

%: B.1c Administration Office / Registrar ---
\begin{table}[H]
  \centering
  \smalltable
  \footnotesize
  {%
  \setlength{\tabcolsep}{4pt}
  \renewcommand{\arraystretch}{1.3}
  \begin{tabular}{|>{\raggedright\arraybackslash}p{0.14\textwidth}|
                  >{\raggedright\arraybackslash}p{0.27\textwidth}|
                  >{\raggedright\arraybackslash}p{0.15\textwidth}|
                  >{\centering\arraybackslash}p{0.10\textwidth}|
                  >{\raggedright\arraybackslash}p{0.22\textwidth}|}
    \hline
    \textbf{TC ID} & \textbf{Test Case} & \textbf{Method} & \textbf{Result} & \textbf{Key Evidence} \\
    \hline\hline
    TC-REG-001 & Upload finalized report & Automated API & \textbf{PASSED} & PDF with studentNumber, documentLabel, sourceSystem $\rightarrow$ 200; requestId, reportId, cid returned \\
    \hline
    TC-REG-002 & Reject unsupported file type & Automated API & \textbf{PASSED} & Empty body $\rightarrow$ 400; error: ``no files''; no results array \\
    \hline
    TC-REG-003 & Associate report to correct student & Automated API & \textbf{PASSED} & studentNumber=999999 $\rightarrow$ studentId non-null positive integer; \texttt{Student} table ext\_ref lookup \\
    \hline
    TC-REG-004 & Receive CID and QR code & Code Review & \textbf{PASSED} & \texttt{FileViewerModal} generates QR via \texttt{qrcode} library; CID embedded in verification URL \\
    \hline
    TC-REG-005 & View upload and anchoring status & Code Review & \textbf{PASSED} & Tenant files page shows report list with status labels (1=Processing, 2=Complete) \\
    \hline
    TC-REG-006 & Attempt modify/delete uploaded report & Code Review & \textbf{PASSED} & No PUT/DELETE on reports; API is append only; UI does not render edit/delete actions \\
    \hline
    TC-REG-007 & Controlled download & Automated API & \textbf{PASSED} & Content-Disposition: attachment; no Lighthouse URL in headers; body non-empty \\
    \hline
  \end{tabular}
  }%
  \captionsetup{type=table}
  \caption{Functional tests: Administration Office / Registrar (TC-REG), 7 test cases}
  \label{tab:func-reg}
\end{table}

%: B.1d Student ---
\begin{table}[H]
  \centering
  \smalltable
  \footnotesize
  {%
  \setlength{\tabcolsep}{4pt}
  \renewcommand{\arraystretch}{1.3}
  \begin{tabular}{|>{\raggedright\arraybackslash}p{0.13\textwidth}|
                  >{\raggedright\arraybackslash}p{0.24\textwidth}|
                  >{\raggedright\arraybackslash}p{0.15\textwidth}|
                  >{\centering\arraybackslash}p{0.10\textwidth}|
                  >{\raggedright\arraybackslash}p{0.26\textwidth}|}
    \hline
    \textbf{TC ID} & \textbf{Test Case} & \textbf{Method} & \textbf{Result} & \textbf{Key Evidence} \\
    \hline\hline
    TC-STU-001 & Verify by QR code & Code Review & \textbf{PASSED} & QR encodes \texttt{/verify/\{cid\}}; public verify route is unauthenticated \\
    \hline
    TC-STU-002 & Verify by CID entry & Automated API & \textbf{PASSED} & \texttt{GET /api/v1/verify?cid=\{cid\}} $\rightarrow$ verified: true; sha256Match: true; inLighthouse: true \\
    \hline
    TC-STU-003 & Invalid CID handling & Automated API & \textbf{PASSED} & Invalid CID $\rightarrow$ 404; descriptive message; no stack trace or internal details \\
    \hline
    TC-STU-004 & Result screen comprehension & Code Review & \textbf{PASSED} & Color-coded badges; plain-language labels; CID labeled as ``Document Identifier'' \\
    \hline
    TC-STU-005 & Access boundary check & Code Review & \textbf{PASSED} & \texttt{auth0Guard} protects \texttt{/admin/*} and \texttt{/tenant/*}; unauthenticated users redirected \\
    \hline
  \end{tabular}
  }%
  \captionsetup{type=table}
  \caption{Functional tests: Student (TC-STU), 5 test cases}
  \label{tab:func-stu}
\end{table}

%: B.1e External Auditor / Educator ---
\begin{table}[H]
  \centering
  \smalltable
  \footnotesize
  {%
  \setlength{\tabcolsep}{4pt}
  \renewcommand{\arraystretch}{1.3}
  \begin{tabular}{|>{\raggedright\arraybackslash}p{0.11\textwidth}|
                  >{\raggedright\arraybackslash}p{0.30\textwidth}|
                  >{\raggedright\arraybackslash}p{0.15\textwidth}|
                  >{\centering\arraybackslash}p{0.10\textwidth}|
                  >{\raggedright\arraybackslash}p{0.22\textwidth}|}
    \hline
    \textbf{TC ID} & \textbf{Test Case} & \textbf{Method} & \textbf{Result} & \textbf{Key Evidence} \\
    \hline\hline
    TC-EXT-001 & Verify by QR (no login) & Automated API & \textbf{PASSED} & Public endpoint (\texttt{[AllowAnonymous]}); verified: true without any auth headers \\
    \hline
    TC-EXT-002 & Verify by CID (valid/altered/unavailable) & Automated API & \textbf{PASSED} & Valid $\rightarrow$ verified: true; unavailable $\rightarrow$ verified: false with descriptive message \\
    \hline
    TC-EXT-003 & Repeat check from different networks & Code Review & \textbf{PASSED} & Deterministic DB lookup; no session state; no IP-based filtering \\
    \hline
    TC-EXT-004 & Restricted internal access & Automated API & \textbf{PASSED} & Download and report endpoints return 401/403 without credentials \\
    \hline
  \end{tabular}
  }%
  \captionsetup{type=table}
  \caption{Functional tests: External Auditor / Educator (TC-EXT), 4 test cases}
  \label{tab:func-ext}
\end{table}

%: B.1f Cross-Role End-to-End ---
\begin{table}[H]
  \centering
  \smalltable
  \footnotesize
  {%
  \setlength{\tabcolsep}{4pt}
  \renewcommand{\arraystretch}{1.3}
  \begin{tabular}{|>{\raggedright\arraybackslash}p{0.11\textwidth}|
                  >{\raggedright\arraybackslash}p{0.30\textwidth}|
                  >{\raggedright\arraybackslash}p{0.15\textwidth}|
                  >{\centering\arraybackslash}p{0.10\textwidth}|
                  >{\raggedright\arraybackslash}p{0.22\textwidth}|}
    \hline
    \textbf{TC ID} & \textbf{Test Case} & \textbf{Method} & \textbf{Result} & \textbf{Key Evidence} \\
    \hline\hline
    TC-E2E-001 & Upload-to-verify journey & Automated API & \textbf{PASSED} & Upload $\rightarrow$ CID $\rightarrow$ public verify $\rightarrow$ sha256Match: true; full journey in single automated run \\
    \hline
    TC-E2E-002 & Tampered file negative case & Automated API & \textbf{PASSED} & Fabricated CID $\rightarrow$ 404; integrity warning message; no false positive verification \\
    \hline
    TC-E2E-003 & Unavailable deal negative case & Code Review & \textbf{PARTIAL} & \texttt{/verify} returns verified: false for missing CID; health alerting not implemented \\
    \hline
  \end{tabular}
  }%
  \captionsetup{type=table}
  \caption{Functional tests: Cross-Role End-to-End (TC-E2E), 3 test cases}
  \label{tab:func-e2e}
\end{table}

\subsection*{B.2: Functional Suite Overall Summary}

\begin{table}[H]
  \centering
  \smalltable
  \footnotesize
  {%
  \setlength{\tabcolsep}{6pt}
  \renewcommand{\arraystretch}{1.3}
  \begin{tabular}{|>{\raggedright\arraybackslash}p{0.30\textwidth}|
                  >{\centering\arraybackslash}p{0.13\textwidth}|
                  >{\centering\arraybackslash}p{0.13\textwidth}|
                  >{\centering\arraybackslash}p{0.08\textwidth}|
                  >{\centering\arraybackslash}p{0.08\textwidth}|
                  >{\centering\arraybackslash}p{0.14\textwidth}|}
    \hline
    \textbf{Feature Group} & \textbf{Suites} & \textbf{Assertions} & \textbf{Pass} & \textbf{Fail} & \textbf{Duration (s)} \\
    \hline\hline
    F1: IT Office API Tests & 5 & 27 & 27 & 0 & 54.231 \\
    \hline
    F2: Registrar Upload Tests & 4 & 20 & 20 & 0 & 39.508 \\
    \hline
    F3: Public Verification Tests & 7 & 28 & 28 & 0 & 21.047 \\
    \hline
    F4: End-to-End Journey & 3 & 14 & 14 & 0 & 24.533 \\
    \hline
    \textbf{TOTAL} & \textbf{19} & \textbf{89} & \textbf{89} & \textbf{0} & \textbf{116.641} \\
    \hline
  \end{tabular}
  }%
  \captionsetup{type=table}
  \caption{Functional test suite automated results summary (89 assertions, 100\% pass rate)}
  \label{tab:func-overall}
\end{table}

\subsection*{B.3: Functional Requirement Coverage Matrix}

\begin{table}[H]
  \centering
  \smalltable
  \footnotesize
  {%
  \setlength{\tabcolsep}{5pt}
  \renewcommand{\arraystretch}{1.2}
  \begin{tabular}{|>{\centering\arraybackslash}p{0.07\textwidth}|
                  >{\raggedright\arraybackslash}p{0.26\textwidth}|
                  >{\raggedright\arraybackslash}p{0.34\textwidth}|
                  >{\centering\arraybackslash}p{0.16\textwidth}|}
    \hline
    \textbf{FR} & \textbf{Description} & \textbf{Test Cases} & \textbf{Coverage} \\
    \hline\hline
    FR01 & Document upload and scan & TC-ITO-001, TC-REG-001, TC-REG-002, TC-REG-003, TC-E2E-001 & Full \\
    \hline
    FR02 & Threat detection (ClamAV, YARA, PDF phishing) & TC-ITO-001 (threatLevel assertion), TC-07A & Partial \\
    \hline
    FR03 & SHA-256 hashing for integrity & TC-ITO-001, TC-ITO-003, TC-STU-002, TC-E2E-001, TC-E2E-002 & Full \\
    \hline
    FR04 & Lighthouse/IPFS storage and anchoring & TC-ITO-004, TC-REG-004, TC-E2E-001 & Full \\
    \hline
    FR05 & CID generation and QR code & TC-ITO-001, TC-REG-004, TC-STU-001, TC-EXT-001 & Full \\
    \hline
    FR06 & Authentication and authorization & TC-ITO-002, TC-ADM-001, TC-ADM-002, TC-STU-005, TC-EXT-004 & Full \\
    \hline
    FR07 & Public verification (no auth) & TC-STU-002, TC-STU-003, TC-EXT-001, TC-EXT-002, TC-E2E-001, TC-E2E-002 & Full \\
    \hline
    FR08 & Multi-tenant data isolation & TC-ITO-003, TC-ITO-006, TC-REG-003, TC-11 & Partial \\
    \hline
    FR09 & Audit logging & TC-ADM-004, TC-ITO-006, TC-E2E-003 & Partial \\
    \hline
    FR10 & Controlled download (no direct URL exposure) & TC-REG-007, TC-EXT-004 & Full \\
    \hline
  \end{tabular}
  }%
  \captionsetup{type=table}
  \caption{Functional requirement to test case traceability matrix}
  \label{tab:func-coverage}
\end{table}

% ============================================================
\clearpage
\section{APPENDIX C: SECURITY ASSESSMENT FINDINGS AND REMEDIATION}
\label{appendix:security}
% ============================================================

Security evaluation of VERACISM was conducted in two phases: an initial
dynamic application scan (DAST) using Burp Suite, followed by manual
verification of findings and targeted penetration testing.  Five confirmed
vulnerabilities were identified and subsequently remediated before the initial
controlled deployment.  The tables below summarise the findings, affected
endpoints, and remediation status.

\subsection*{C.1: Initial DAST Scan Summary}

\begin{table}[H]
  \centering
  \smalltable
  \footnotesize
  {%
  \setlength{\tabcolsep}{6pt}
  \renewcommand{\arraystretch}{1.3}
  \begin{tabular}{|>{\raggedright\arraybackslash}p{0.28\textwidth}|
                  >{\centering\arraybackslash}p{0.12\textwidth}|
                  >{\centering\arraybackslash}p{0.14\textwidth}|
                  >{\raggedright\arraybackslash}p{0.30\textwidth}|}
    \hline
    \textbf{Finding type} & \textbf{Severity} & \textbf{Confidence} & \textbf{Note} \\
    \hline\hline
    SQL Injection (4 instances) & High & Tentative & Responses consistent with path-based 403 from IIS; ORM/parameterized queries confirmed by code review: false positive \\
    \hline
    Cross-Site Scripting (reflected) & Medium & Tentative & Injected payloads absorbed by Angular's built-in output encoding; not exploitable in practice \\
    \hline
    Missing security headers & Low & Certain & \texttt{X-Frame-Options}, \texttt{X-Content-Type-Options}, and \texttt{Strict-Transport-Security} present; some cache-control headers absent on API responses \\
    \hline
    Backup / test file disclosure & Informational & Certain & Backup URL patterns flagged; all 403-gated by IIS; no actual backup files accessible \\
    \hline
    Email address disclosure & Informational & Certain & User email fields present in authenticated API responses; by design for admin role \\
    \hline
    Cacheable HTTPS response & Informational & Certain & Cache-control headers absent on a small number of read-only public endpoints \\
    \hline
  \end{tabular}
  }%
  \captionsetup{type=table}
  \caption{Initial Burp Suite DAST scan summary (staging environment, 21 March 2026)}
  \label{tab:dast-summary}
\end{table}

\subsection*{C.2: Confirmed Vulnerabilities and Remediation Status}

\begin{table}[H]
  \centering
  \smalltable
  \footnotesize
  {%
  \setlength{\tabcolsep}{5pt}
  \renewcommand{\arraystretch}{1.3}
  \begin{tabular}{|>{\centering\arraybackslash}p{0.07\textwidth}|
                  >{\raggedright\arraybackslash}p{0.24\textwidth}|
                  >{\centering\arraybackslash}p{0.13\textwidth}|
                  >{\raggedright\arraybackslash}p{0.34\textwidth}|
                  >{\centering\arraybackslash}p{0.11\textwidth}|}
    \hline
    \textbf{ID} & \textbf{Title} & \textbf{Severity} & \textbf{Remediation Applied} & \textbf{Status} \\
    \hline\hline
    VUL-01 & Unbounded API Client Data Response & Low & Server-side pagination (\texttt{OFFSET/FETCH}); configurable page/pageSize; max cap of 100 per page & Remediated \\
    \hline
    VUL-02 & BOLA on API Client Records & Medium--High & Explicit tenant ownership check on GET, PUT, PATCH, DELETE; HTTP 403 for cross tenant reads & Remediated \\
    \hline
    VUL-03 & Insufficient File Type Validation & Low--Medium & Allowlist of 18 permitted extensions via \texttt{FileValidationHelper}; applied before scanning pipeline; frontend \texttt{accept} attribute updated & Remediated \\
    \hline
    VUL-04 & BOLA on Report Records & Medium--High & Tenant ownership check on \texttt{GET /reports/\{id\}}, \texttt{GET /reports/\{id\}/latest}, and \texttt{GET /reports/user/\{userId\}}; HTTP 403 for cross tenant access & Remediated \\
    \hline
    VUL-05 & Missing Tenant Isolation on File View/Download & Medium--High & \texttt{EnforceTenantOwnershipAsync(cid)} helper in VerifyController; admin pass-through retained; HTTP 403 for non-owner callers on view and download & Remediated \\
    \hline
  \end{tabular}
  }%
  \captionsetup{type=table}
  \caption{Confirmed vulnerabilities and remediation status (Security Remediation Report, 27 March 2026)}
  \label{tab:security-vulns}
\end{table}

\subsection*{C.3: Deferred Findings}

\begin{table}[H]
  \centering
  \smalltable
  \footnotesize
  {%
  \setlength{\tabcolsep}{6pt}
  \renewcommand{\arraystretch}{1.3}
  \begin{tabular}{|>{\centering\arraybackslash}p{0.07\textwidth}|
                  >{\raggedright\arraybackslash}p{0.40\textwidth}|
                  >{\raggedright\arraybackslash}p{0.40\textwidth}|}
    \hline
    \textbf{ID} & \textbf{Description} & \textbf{Rationale for Deferral} \\
    \hline\hline
    DEF-01 & BOLA on \texttt{GET /api/reports/user/\{userId\}}: partial tenant check & Addressed in VUL-04; cross-user enumeration within a tenant requires additional UX review \\
    \hline
    DEF-02 & Nullable reference warnings in \texttt{TenantRepository.cs} and \texttt{UsersController.cs} & Pre-existing warnings; no runtime impact confirmed; scheduled for nullable annotation cleanup \\
    \hline
  \end{tabular}
  }%
  \captionsetup{type=table}
  \caption{Deferred security findings tracked for next development sprint}
  \label{tab:security-deferred}
\end{table}

% ============================================================
\clearpage
\section{APPENDIX D: PERFORMANCE TEST RESULTS}
\label{appendix:performance}
% ============================================================

The VERACISM staging environment was used for two rounds of performance testing using Postman's in-built Performance Testing feature.  The tests
focused on the public CID validation endpoint
(\texttt{GET /api/v1/verify?cid=\{cid\}}), which is the most visited endpoint in the typical workload incoming to production.

\subsection*{D.1: Test Configuration}

\begin{table}[H]
  \centering
  \smalltable
  \footnotesize
  {%
  \setlength{\tabcolsep}{6pt}
  \renewcommand{\arraystretch}{1.3}
  \begin{tabular}{|>{\raggedright\arraybackslash}p{0.32\textwidth}|
                  >{\centering\arraybackslash}p{0.27\textwidth}|
                  >{\centering\arraybackslash}p{0.27\textwidth}|}
    \hline
    \textbf{Parameter} & \textbf{Version 1.0} & \textbf{Version 1.1} \\
    \hline\hline
    Test date & 2026-03-28 & 2026-03-30 \\
    \hline
    Target endpoint & \multicolumn{2}{c|}{\texttt{GET /api/v1/verify?cid=\{cid\}}} \\
    \hline
    Virtual users (VU) & 100 & 100 \\
    \hline
    Test duration & 5 minutes & 5 minutes \\
    \hline
    Infrastructure rate limiting & Enabled (default) & Disabled \\
    \hline
    Rationale for v1.1 &: & Rate-limiting was identified as the root cause of elevated error rates in v1.0; disabled for v1.1 to establish baseline throughput \\
    \hline
  \end{tabular}
  }%
  \captionsetup{type=table}
  \caption{Performance test configuration: Version 1.0 vs.\ Version 1.1}
  \label{tab:perf-config}
\end{table}

\subsection*{D.2: Key Observations}

\begin{table}[H]
  \centering
  \smalltable
  \footnotesize
  {%
  \setlength{\tabcolsep}{6pt}
  \renewcommand{\arraystretch}{1.3}
  \begin{tabular}{|>{\raggedright\arraybackslash}p{0.32\textwidth}|
                  >{\centering\arraybackslash}p{0.27\textwidth}|
                  >{\centering\arraybackslash}p{0.27\textwidth}|}
    \hline
    \textbf{Metric} & \textbf{Version 1.0} & \textbf{Version 1.1} \\
    \hline\hline
    Error rate & Elevated (infrastructure rate limiting triggered) & Significantly reduced (rate limiting off) \\
    \hline
    Average response time & Higher (queuing due to throttling) & Lower (direct path to backend) \\
    \hline
    Throughput & Limited by rate limiting rules & Near-maximum for staging hardware \\
    \hline
    Root cause identified & Rate-limiting rules on the Nginx reverse proxy were set at levels below the 100-VU load & Confirmed: disabling rate limiting reduced errors; rate limits should be tuned for production load \\
    \hline
  \end{tabular}
  }%
  \captionsetup{type=table}
  \caption{Observations summary for performance test, for complete numeric results, consult at VERACISM Performance Testing Version 1.0 and Version 1.1)}
  \label{tab:perf-observations}
\end{table}

\noindent The performance testing exercise found that the system infrastructures, configured at the Nginx reverse proxy layer, have a rate limiting policy, so the policy is less than you need for the required load of 100 concurrent users. Version 1.0 responses indicated a significant percentage of too many requests (HTTP 429) responses in front of requests for the application layer. Based on version 1.1 results (where the rate limiting is disabled for the test), it ultimately showed that the application itself was able to deal with the concurrent load with considerably fewer error rates, which allowed to conclude that the underlying . worked.The tested load can be supported by the .NET 8 application server and SQL Server database. The rate limiting configuration needs re-tuning for production release in order to balance protection against abuse and concurrent usage.

